\chapter{Sleep-Related Rhythmic Movement Disorder}
\index{Sleep-Related Rhythmic Movement Disorder}
\label{sec:Sleep-Related Rhythmic Movement Disorder}

\section{Overview}
Sleep-related rhythmic movement disorder (RMD) is a sleep-related movement disorder involving stereotypical, repetitive movements of large muscle groups occurring primarily during the transition from wakefulness to sleep or during NREM sleep. The most common forms are head banging, head rolling, and body rocking. Prevalence is up to 60\% in infants (physiologic), declining to less than 1\% in adults. This condition is among the most common mental health disorders worldwide. It affects people across all age groups, socioeconomic backgrounds, and geographic regions. This condition affects a significant number of people worldwide and can range from mild to severe in presentation. Early recognition and appropriate management are essential for optimal outcomes. A thorough understanding of the underlying mechanisms guides treatment decisions. Patient education and self-management play important roles in long-term care.
\section{Causes}
The underlying mechanisms involve complex interactions between genetic predisposition and environmental factors. Disruption of normal physiological processes leads to the characteristic features of the condition. Multiple contributing factors may act synergistically to trigger or worsen the condition.

\begin{itemize}[noitemsep]
  \item This condition happens because of unknown Disease mechanism
  \item In adults, it is linked to neurodevelopmental disorders, anxiety, and OSA.
\end{itemize}
\section{Risk Factors}
Genetic predisposition and family history Advancing age Lifestyle factors (diet, exercise, smoking, alcohol) Environmental exposures Underlying medical conditions Medication use Stress and psychological factors Socioeconomic factors

\begin{itemize}[noitemsep]
  \item Genetic predisposition and family history
  \item Advancing age
  \item Lifestyle factors (diet, exercise, smoking, alcohol)
  \item Environmental exposures
  \item Underlying medical conditions
  \item Medication use
\end{itemize}
\section{Signs and Symptoms}
You may experience rhythmic movements at a frequency of 0.5--2 Hz lasting less than 15 minutes. Symptoms vary depending on the severity and extent of the condition Pain or discomfort in the affected area Functional impairment affecting daily activities Changes in appearance or function of affected tissues Systemic symptoms may include fatigue, fever, or weight changes Symptoms may develop gradually or appear suddenly

\begin{itemize}[noitemsep]
  \item Pain or discomfort in the affected area
  \item Functional impairment affecting daily activities
  \item Changes in appearance or function of affected tissues
  \item Systemic symptoms may include fatigue, fever, or weight changes
  \item Symptoms may develop gradually or appear suddenly
\end{itemize}
\section{Progression and Stages}
The condition may progress through different stages of severity over time. Early stages often involve mild or subtle symptoms that may be overlooked. Moderate stages involve more noticeable symptoms affecting daily function. Advanced stages may involve significant impairment and complications.
\section{Complications}
\begin{itemize}[noitemsep]
  \item Disease progression leading to worsening symptoms
  \item Secondary infections or injuries
  \item Side effects of treatment
  \item Reduced quality of life
  \item Emotional and psychological impact
\end{itemize}
\section{Diagnosis}
\begin{itemize}[noitemsep]
  \item Doctors usually diagnose this based on your symptoms and a physical exam
  \item Polysomnography with synchronized video recording is indicated when injury is suspected or to rule out seizure.
  \item Comprehensive medical history and physical examination
  \item Laboratory tests to assess organ function
  \item Imaging studies to visualize affected structures
  \item Specialized testing depending on the affected system
  \item Consultation with relevant specialists
\end{itemize}
\section{Prevention}
\begin{itemize}[noitemsep]
  \item Maintain a healthy lifestyle with balanced diet and exercise
  \item Avoid known risk factors
  \item Attend regular medical check-ups for early detection
  \item Follow recommended screening guidelines
\end{itemize}
\section{Treatment Options}
\begin{itemize}[noitemsep]
  \item Treatment typically involves primarily reassurance and safety measures
  \item Clonazepam, hypnotics, or behavioral therapy may be used in persistent or injurious adult cases.
  \item Medications to manage symptoms and address underlying mechanisms
  \item Lifestyle modifications and self-care measures
  \item Physical therapy and rehabilitation
  \item Surgical intervention when indicated
  \item Regular monitoring and follow-up care
\end{itemize}
\section{Living with It}
\begin{itemize}[noitemsep]
  \item Follow your treatment plan as prescribed
  \item Maintain regular follow-up with your healthcare provider
  \item Make necessary lifestyle adjustments
  \item Seek support from family, friends, or support groups
  \item Stay informed about your condition
\end{itemize}
\section{Outlook}
\begin{itemize}[noitemsep]
  \item Outlook is non-cancerous
  \item Most children outgrow the condition.
\end{itemize}
